%!TEX root = ../hutbthesis_main.tex

\chapter{仿真系统构建与实现}

\section{系统总体架构与软硬件环境}

仿真平台由五个层次组成：Ubuntu 18.04 提供系统运行基础；ROS Noetic
提供节点通信、参数服务器、TF 与消息系统；Gazebo 11
提供刚体动力学、碰撞、接触与传感器仿真；gazebo\_ros\_pkgs 和 gazebo\_ros\_control 将
Gazebo 的模型、传感器与控制接口映射到
ROS；项目自定义节点完成飞机辅助感知、局部规划、牵引对接控制与任务调度\textsuperscript{[11-13,16]}。Gazebo
官方 ROS 教程明确指出，gazebo\_ros\_pkgs 负责 ROS 与 Gazebo
的接口层，而 gazebo\_ros\_control 通过解析 URDF
中的 \textless transmission\textgreater 与插件配置，实现对仿真关节的
controllers 接入\textsuperscript{[13,16]}。

在实际工程目录上，采用四包式组织方式：tug\_description 保存
xacro/URDF、网格模型与材料；tug\_gazebo 保存 world、spawn
与传感器插件；tug\_control 保存控制器参数、接口桥接与控制节点；tug\_bringup 保存联调
launch 文件。

Gazebo
的 roslaunch 教程还给出了较为稳定的包层级约定，即 MYROBOT\_description/urdf、MYROBOT\_gazebo/worlds 与 MYROBOT\_gazebo/launch 的分离组织，这一约定适合本课题的工程化管理\textsuperscript{[11-16]}。

下表给出本文采用的基准软件栈。（表4-1）

\begin{center}\textbf{表4-1 基准软件栈}\end{center}

\begin{longtable}[]{@{}
  >{\raggedright\arraybackslash}p{(\columnwidth - 4\tabcolsep) * \real{0.2678}}
  >{\raggedright\arraybackslash}p{(\columnwidth - 4\tabcolsep) * \real{0.3713}}
  >{\raggedright\arraybackslash}p{(\columnwidth - 4\tabcolsep) * \real{0.3608}}@{}}
\toprule\noalign{}
\begin{minipage}[b]{\linewidth}\raggedright
组件
\end{minipage} & \begin{minipage}[b]{\linewidth}\raggedright
版本/状态
\end{minipage} & \begin{minipage}[b]{\linewidth}\raggedright
用途
\end{minipage} \\
\midrule\noalign{}
\endhead
\bottomrule\noalign{}
\endlastfoot
Ubuntu & 18.04 LTS & 主机操作系统 \\
ROS & Noetic Ninjemys & 节点通信、参数、消息、TF \\
Gazebo & 11 & 刚体动力学、传感器与场景仿真 \\
gazebo\_ros\_pkgs & Noetic 对应发行 & Gazebo 与 ROS 接口桥 \\
gazebo\_ros\_control & Noetic 对应发行 & ros\_control 接入 Gazebo \\
ros\_control & ROS 1 经典控制框架 & 控制器抽象与管理 \\
robot\_state\_publisher & Noetic 对应发行 & 根据 URDF
+ joint\_states发布 TF \\
tf2\_ros & Noetic 对应发行 & 坐标变换管理 \\
xacro & Noetic 对应发行 & 宏化 URDF 生成 \\
image\_pipeline & Noetic 对应发行 & 相机标定与图像处理链 \\
OpenCV & ROS Noetic 兼容版本 & 飞机辅助感知标识识别 \\
\end{longtable}

如下系统架构图适合在论文中绘制为 PNG 或 SVG，用于说明 ROS/Gazebo
联合仿真的模块层级关系。其图元关系与 Gazebo ROS
插件、ros\_control 和 robot\_state\_publisher 的官方工作方式一致。

\section{牵引车模型与传感器建模}

URDF 从结构上将机器人表示为 link--joint 树；Gazebo
官方教程进一步指出，若要把 URDF
真正用于动力学仿真，除 visual 外还必须补充 collision、inertial、Gazebo
专用 \textless gazebo\textgreater 标签以及 ros\_control 所需的 \textless transmission\textgreater{}\textsuperscript{[13-16]}。同时，URDF
本身无法描述世界中的模型摆放、闭环关节与部分接触参数，因此世界摆放与环境物体建议放在 .world 中，机器人本体保持单一树形结构\textsuperscript{[12-14]}。

对于双桥转向牵引车，面向第三章算法实现，推荐将车体模型拆分为以下
link：、base\_link、front\_ateer、back\_steer、四个车轮 link、相机
link、深度相机 link、激光雷达 link，以及飞机辅助感知所需的可视化或调试
link。关节层面至少需要：前桥转向关节、后桥转向关节、四个轮子的转动关节，以及各传感器与车体的固定关节。该拆分方式既符合
URDF
树结构，也便于单独对前后桥转向与轮速控制器建模\textsuperscript{[14,16]}。

车辆基础仿真与机械模型都基于SOLIDWORKS实现。考虑到ROS与Gazebo在虚拟机环境下性能受限的影响，将模型分为两种建立------仿真模型与机械模型。其中仿真面向基本运行逻辑的测试与较为极端的算法测试，保证项目推进速度的同时最大程度降低硬件成本。其结构较为简单，隐藏了不必要的机械结构部分，只保留最基础的车辆外形与少量运动模块；机械模型是最终提交到现实生产的版本，其结构更为复杂，相比仿真模型要考虑更多的机械及电气布局。首篇幅限制，下面只给出仿真模型（图4-1）

\begin{figure}[htbp]
    \centering
    \includegraphics[width=0.82\linewidth]{image8.png}
    \caption{模型图总览}
\end{figure}

模型设计完成后，需要对各个活动关节进行定义与功能区分。根据实际功能------continuous（关节持续转动）、fixed（关节固定）、revolute（关节有限制有角度旋转），对车轮、转向桥以及功能部件定义。下表给出面向论文复现的一组牵引车
URDF
结构定义。表中“用途”与“建议参数”基于本项目双桥转向、小尺度机场地面车辆仿真的工程需求给出；若项目后续补充了
CAD/BOM 或实测参数，应以实测量替换。（表4-2）

\begin{center}\textbf{表4-2 车辆主要URDF结构定义}\end{center}

\begin{longtable}[]{@{}
  >{\raggedright\arraybackslash}p{(\columnwidth - 6\tabcolsep) * \real{0.2552}}
  >{\raggedright\arraybackslash}p{(\columnwidth - 6\tabcolsep) * \real{0.2294}}
  >{\raggedright\arraybackslash}p{(\columnwidth - 6\tabcolsep) * \real{0.1982}}
  >{\raggedright\arraybackslash}p{(\columnwidth - 6\tabcolsep) * \real{0.3172}}@{}}
\toprule\noalign{}
\begin{minipage}[b]{\linewidth}\raggedright
组件
\end{minipage} & \begin{minipage}[b]{\linewidth}\raggedright
类型
\end{minipage} & \begin{minipage}[b]{\linewidth}\raggedright
命名
\end{minipage} & \begin{minipage}[b]{\linewidth}\raggedright
说明
\end{minipage} \\
\midrule\noalign{}
\endhead
\bottomrule\noalign{}
\endlastfoot
车体根坐标 & link & base\_link & 仅作平面参考，不带惯量 \\
主车体 & link & base\_link & 车体惯量、碰撞主体 \\
前桥 & link & front\_steer\_link & 前转向桥中间体 \\
后桥 & link & back\_steer\_link & 后转向桥中间体 \\
前桥转向 & revolute & front\_steer & 绕 z 轴转向 \\
后桥转向 & revolute & rear\_steer & 绕 z 轴转向 \\
相机 & fixed & front\_camera\_joint & 识别飞机辅助码主相机 \\
深度相机 & fixed & depth\_camera\_joint & 近距深度感知 \\
激光雷达 & fixed & lidar\_joint & 前向近距安全监测 \\
\end{longtable}

在此示例部分关节定义：

  \textless link

    name="base\_link"\textgreater{}

    \textless inertial\textgreater{}

      \textless origin

        xyz="0.052294 1.5562E-06 0.061913"

        rpy="0 0 0" /\textgreater{}

      \textless mass

        value="34.871" /\textgreater{}

      \textless inertia

        ixx="0.51812"

        ixy="-2.2269E-05"

        ixz="0.18429"

        iyy="2.6761"

        iyz="-2.8808E-06"

        izz="3.0257" /\textgreater{}

    \textless/inertial\textgreater{}

    \textless visual\textgreater{}

      \textless origin

        xyz="0 0 0"

        rpy="0 0 0" /\textgreater{}

      \textless geometry\textgreater{}

        \textless mesh

          filename="package://zongzhuang/meshes/base\_link.STL" /\textgreater{}

      \textless/geometry\textgreater{}

      \textless material

        name=""\textgreater{}

        \textless color

          rgba="0.79216 0.81961 0.93333 1" /\textgreater{}

      \textless/material\textgreater{}

    \textless/visual\textgreater{}

    \textless collision\textgreater{}

      \textless origin

        xyz="0 0 0"

        rpy="0 0 0" /\textgreater{}

      \textless geometry\textgreater{}

        \textless mesh

          filename="package://zongzhuang/meshes/base\_link.STL" /\textgreater{}

      \textless/geometry\textgreater{}

    \textless/collision\textgreater{}

  \textless/link\textgreater{}

\section{任务与模式管理}

任务与模式管理器在整个控制系统中起着核心作用\textsuperscript{[25,28]}。该模块的主要任务是根据感知系统提供的信息（如目标身份确认、位姿识别结果、障碍物检测等），决定当前的控制模式。具体模式包括\textsuperscript{[5,7,28]}：

搜索模式：当目标身份未确认时，系统处于搜索模式，牵引车保持静止或缓慢扫描。

自动接近模式：当目标身份确认并且距离较远时，系统启动自动接近模式，牵引车沿预设路径接近目标飞机。

姿态对准模式：当牵引车距离目标飞机较近时，系统切换到姿态对准模式，通过调整转向角度和速度逐步减少横向误差和姿态误差。

微动对接模式：在接近阶段，系统切换到微动对接模式，牵引车以非常低的速度精确接近并最终完成对接。

停车模式：当目标对接完成或无法继续推进时，系统进入停车模式，发布停车命令。

\section{系统全流程测试}

为验证自动牵引对接系统在完整作业流程中的可行性与连续性，本文在ROS与Gazebo联合仿真环境中设计了自动牵引车全流程测试。测试过程按照机场地面牵引作业的基本逻辑展开，主要包括车辆前往飞行器区域、车辆对接飞行器、车辆推出飞行器、解除牵引以及车辆撤离五个阶段\textsuperscript{[5-7]}。通过该测试，可以较为直观地观察牵引车在不同作业阶段下的运动状态、姿态变化以及控制指令执行效果。

1. 车辆前往飞行器区域

测试开始，牵引车从模拟航站楼附近出发，根据地勤车辆道路向飞行器所在区域移动。在该阶段，系统主要完成车辆启动、路径跟踪和环境感知等任务。牵引车通过传感器获取周围环境信息，并结合控制算法调整自身速度与转向角，使车辆能够沿规划路径稳定行驶至飞行器附近。

在仿真过程中，车辆需要保持较低且平稳的行驶速度，避免因速度过高造成轨迹偏差或姿态调整不及时。当前往飞行器区域时，系统重点关注车辆行驶方向是否稳定、路径跟踪是否连续，以及车辆是否能够准确进入后续对接准备区域。（图4-2）

\begin{figure}[htbp]
    \centering
    \includegraphics[width=0.82\linewidth]{image9.png}
    \caption{车辆在航站楼下运行，前往目标飞行器}
\end{figure}

2. 车辆对接飞行器

当牵引车进入飞行器附近区域后，系统切换至对接控制阶段。该阶段的核心任务是使牵引车与飞行器牵引点逐步对准，并在满足位置和姿态要求的情况下完成接近动作。系统根据感知模块提供的飞行器相对位姿信息，计算牵引车与目标对接点之间的位置误差和航向误差，并通过控制算法修正车辆运动状态。（图4-3）

在对接过程中，牵引车首先进入预对接位置，然后以低速向最终对接点靠近。相比普通路径跟踪阶段，对接阶段对控制精度和车辆稳定性要求更高，因此需要限制车辆速度，并使转向角变化更加平缓，避免出现车身摆动过大或对接偏差过大的情况。（图4-4）

\begin{figure}[htbp]
    \centering
    \includegraphics[width=0.82\linewidth]{image10.png}
    \caption{尝试对接。此时车辆没有正对飞行器，正在通过蟹行平移进行调整}
\end{figure}

\begin{figure}[htbp]
    \centering
    \includegraphics[width=0.82\linewidth]{image11.png}
    \caption{对接成功，拖把末端固定装置放下。等待推出命令}
\end{figure}

该阶段主要验证牵引车是否能够根据飞行器目标位置完成姿态调整，并稳定接近指定对接点。

3. 车辆推出飞行器

完成对接后，系统进入飞行器推出阶段。此时牵引车与飞行器被视为一个整体运动系统，牵引车根据预设推出路径向指定方向行驶，将飞行器从停放区域推出至目标位置。在该阶段，牵引车需要保持稳定牵引状态，并避免转向过大导致飞行器姿态变化不平顺。

推出过程中，控制系统重点保证车辆速度平稳、转向连续，并使飞行器整体运动轨迹满足预设作业要求。由于推出阶段涉及车辆与飞行器之间的相对约束关系，因此相比单车行驶更需要关注运动过程的稳定性和路径偏差。（图4-5------图4-7）

\begin{figure}[htbp]
    \centering
    \includegraphics[width=0.82\linewidth]{image12.png}
    \caption{飞行器被推出过程}
\end{figure}

\begin{figure}[htbp]
    \centering
    \includegraphics[width=0.82\linewidth]{image13.png}
    \caption{飞行器被推出过程}
\end{figure}

\begin{figure}[htbp]
    \centering
    \includegraphics[width=0.82\linewidth]{image14.png}
    \caption{飞行器被推出过程}
\end{figure}

4. 解除牵引

当飞行器被推出至指定位置后，牵引车停止运动，并执行解除牵引动作。该阶段主要用于模拟实际机场作业中牵引车与飞行器之间的脱离过程。系统首先控制车辆减速至停止状态，随后解除与飞行器之间的连接关系，使飞行器保持在目标位置，牵引车进入独立运动状态。

解除牵引阶段的关键在于保证车辆停止平稳，避免因惯性造成飞行器位置偏移。同时，解除牵引后应确认飞行器姿态稳定，牵引车与飞行器之间保持安全距离，为后续车辆撤离提供空间。（图4-6）

\begin{figure}[htbp]
    \centering
    \includegraphics[width=0.82\linewidth]{image15.png}
    \caption{拖把解除，车辆倒退一部分距离等待飞行器报告}
\end{figure}

5. 车辆撤离

完成解除牵引后，牵引车按照预设撤离路径离开飞行器区域。撤离阶段主要验证牵引车在完成作业任务后能否安全驶离，并避免与飞行器或周围环境发生干涉。系统根据撤离方向控制车辆转向与速度，使其逐步远离飞行器，并回到指定待命区域或安全区域。

在该阶段，车辆应保持稳定行驶状态，转向过程应尽量平缓，避免因急转或急停影响仿真效果。车辆撤离后，整个自动牵引对接与推出流程结束。（图4-7）

\begin{figure}[htbp]
    \centering
    \includegraphics[width=0.82\linewidth]{image16.png}
    \caption{车辆撤离}
\end{figure}

\section{本文方法与人工操作以及其它方法对比分析}

为进一步说明本文所设计系统的特点，有必要将其与传统人工牵引方式、固定路径控制方式以及常见自动驾驶控制算法进行对比分析。

\subsection{与传统人工牵引方式的比较}

传统机场牵引作业主要依靠驾驶员、地面指挥员以及现场协同人员完成。其优势在于人工判断灵活，能够根据飞机位置、现场障碍物和突发情况快速调整操作策略，因此在真实复杂环境中的适应性较强。但是，该方式也存在对人员经验依赖较大、操作标准不易统一、重复作业精度受人为因素影响等问题。

本文所提出的方法主要面向仿真环境下的自动化流程验证。相比人工牵引方式，本文方法的优势并不在于真实复杂环境中的经验判断能力，而在于能够将牵引作业过程转化为可建模、可重复、可分析的控制流程。通过ROS节点通信、Gazebo物理仿真、车辆运动控制和目标位姿误差反馈，可以较为清晰地表达车辆接近、对准、推出和撤离等动作过程。

从本项目实际测试过程来看，车辆模型在Gazebo中能够完成基本行驶、转向、刹停和回中等控制动作，并可通过脚本或控制节点实现分阶段作业流程。这说明本文方法在标准化仿真验证方面具有一定价值。下图给出了本文方法与其它方法的比较柱状图（图4-8）。

\begin{figure}[htbp]
    \centering
    \includegraphics[width=0.82\linewidth]{image17.png}
    \caption{本文方法与主流方法以及理论方法比较}
\end{figure}

\subsection{与固定路径或开环脚本控制方式的比较}

相比与其它方式而言，最简单实现机场自动牵引驳接的方式是采用固定路径或开环脚本控制，即预先设定车辆的速度、转向角和运动时间，使车辆按照固定流程完成前进、转弯、倒车或停止等动作。但固定路径或开环控制的局限性也比较明显。由于其控制过程主要依赖预设时间和固定参数，当车辆初始位置、飞机停放位置或仿真环境发生变化时，车辆无法根据实际偏差自动修正运动状态。例如，在本项目调试过程中，车辆曾出现初始缓慢滑动、转向不完全回中、不同速度下运动效果不一致等现象，如果完全依靠固定脚本控制，容易造成路径偏差累积，影响后续对接效果。

本文方法在固定流程的基础上引入了相对位姿误差思想，将车辆与目标对接点之间的位置误差、横向误差和航向误差作为控制依据，并结合预对接点、最终对接点和分阶段模式切换来组织整个作业流程。因此，相较于单纯开环脚本，本文方法具有更强的闭环控制特征，能够更合理地描述自动牵引车在接近目标过程中的修正行为（图4-9）。

\begin{figure}[htbp]
    \centering
    \includegraphics[width=0.82\linewidth]{image18.png}
    \caption{几种方法对接过程中的横向误差收敛对比}
\end{figure}

\subsection{与常见路径跟踪算法的比较}

在移动机器人和智能车辆研究中，常见路径跟踪方法包括PID控制、Pure
Pursuit算法、Stanley算法以及模型预测控制等。

PID控制结构简单、实现方便，适合速度控制、角度控制等单一变量调节。在本项目中，车辆的车轮速度控制器和转向关节控制器本身就需要依赖底层控制参数实现稳定执行。但如果仅采用简单PID完成整体对接任务，则难以同时处理车辆纵向距离、横向偏差和航向角误差之间的耦合关系。牵引车对接飞行器时不仅要“走到目标附近”，还要保证车身姿态与牵引点方向基本一致，因此单一PID控制难以完整表达对接过程。

Pure
Pursuit算法适合一般路径跟踪，能够根据前视点生成较平滑的转向控制量。但在近距离精确对接阶段，车辆与目标点之间距离较短，速度较低，对末端姿态要求较高，单纯依靠前视点跟踪可能会出现末端对准不够精细的问题。本文方法更加关注车辆与对接目标之间的相对位姿关系，因此更适合描述低速、近距离、姿态约束较强的对接场景。

Stanley算法在自动驾驶路径跟踪中应用较多，尤其适合处理中高速行驶过程中的横向误差修正。但机场牵引车对接飞行器属于低速精细作业，车辆运动速度较低，控制重点不是高速稳定循迹，而是小范围内的位置修正和姿态调整。因此，本文采用基于相对误差的分阶段控制思路，比直接套用Stanley方法更符合本课题场景。

MPC方法可以综合考虑车辆运动约束、转向角限制、速度限制和未来轨迹优化，理论上更适合复杂约束下的车辆控制。但MPC对车辆模型精度、求解器实时性和参数调试要求较高。结合本项目实际情况，课题时间有限，重点又在于完成ROS-Gazebo联合仿真系统搭建和全流程测试，直接采用MPC会显著增加实现难度。因此，本文方法没有追求复杂最优控制，而是选择更容易落地的相对位姿误差控制与规则化流程管理，在可实现性和论文工作量之间取得平衡（图4-10）。

\begin{figure}[htbp]
    \centering
    \includegraphics[width=0.82\linewidth]{image19.png}
    \caption{多种方法对接平稳度对比}
\end{figure}

\subsection{对比结果总结}

综合上述分析，本文方法的主要特点并不是在所有指标上优于人工或其他算法，而是在特定研究阶段和特定应用目标下具有较好的适用性。其优势主要体现在以下几个方面：

第一，本文方法能够把机场牵引作业拆解为多个可控阶段，使“前往飞行器区域、对接飞行器、推出飞行器、解除牵引、车辆撤离”形成较完整的仿真流程。

第二，本文基于ROS与Gazebo构建了可重复测试的仿真平台，便于对车辆模型、控制策略和作业流程进行验证，相比直接进行真实车辆测试，风险和成本更低。

第三，本文方法结合车辆相对位姿误差进行控制分析，比单纯开环脚本更能体现自动对接过程中的误差修正思想。

第四，本文采用较为清晰的分层结构，将感知、任务管理、局部目标生成、误差控制和执行层连接起来，便于后续继续扩展传感器融合、目标识别和路径优化算法。

下图能够直观展现各个方法的能效性能对比图（图4-11）。

总而言之，本文算法思路建立在对以上多种手动、自动牵引作业优缺点综合考虑下，最大可能集成上述方法的优点减少缺陷，提升车辆在复杂环境下的自适应能力与任务鲁棒性。

\begin{figure}[htbp]
    \centering
    \includegraphics[width=0.82\linewidth]{image20.png}
    \caption{不同方法的综合能力对比}
\end{figure}

\section{本章小结}

本章围绕ROS-Gazebo联合仿真系统的构建与实现展开，重点说明了仿真软件栈、牵引车模型、传感器挂载、任务模式管理以及ROS与Gazebo之间的接口关系。通过URDF模型、Gazebo场景、ros\_control控制器和ROS节点通信，牵引车模型能够完成加载、转向回中、轮速控制、目标接近和停车保持等基础动作，为第三章提出的感知与控制方法提供了验证环境。后续可在此基础上替换车辆模型、调整控制器参数或增加更复杂的机场地图与传感器配置。为接下来的实际任务部署提供一个快速迭代的实验平台。
